\chapter*{Tóm tắt}
\addcontentsline{toc}{chapter}{Tóm tắt}

Sự phát triển nhanh chóng của các hệ thống mạng viễn thông hiện đại đã đặt ra nhiều thách thức lớn trong quá trình vận hành và xử lý sự cố. Các phương pháp chẩn đoán truyền thống thường dựa vào luật thủ công hoặc các mô hình AI tĩnh, gặp nhiều khó khăn trong việc thích ứng với các lỗi mới hoặc sự thay đổi của cấu trúc mạng. Nhằm giải quyết vấn đề này, đề tài tập trung nghiên cứu, xây dựng và đánh giá hệ thống **Self-Evolving AI Agents** hỗ trợ tự động xử lý sự cố mạng dựa trên nền tảng **NIKA** (A Network Arena for Benchmarking AI Agents on Network Troubleshooting).

Nền tảng NIKA cung cấp môi trường mô phỏng mạng thực tế dựa trên Kathará, hỗ trợ giả lập nhiều sự cố phức tạp (như lỗi cấu hình định tuyến BGP/OSPF, lỗi thiết bị mạng, SDN controller và P4 data-plane) cùng cơ chế đánh giá hiệu năng chẩn đoán tự động. Nghiên cứu này hiện thực hóa hai mô-đun cốt lõi giúp AI Agent có khả năng tự tiến hóa:
\begin{enumerate}
    \item **Mô-đun Bộ Nhớ Quy Trình Tiến Hóa (Composable Procedural Memory Module):** Sử dụng lưu trữ kết hợp PostgreSQL và Qdrant để lưu trữ kinh nghiệm chẩn đoán dưới dạng các ghi chú nguyên tử (atomic notes). Mô-đun sử dụng cơ chế nén vết chẩn đoán (trajectory compaction) lấy cảm hứng từ LightMem, truy xuất theo thuộc tính ngữ cảnh (attribute-aware retrieval) lấy cảm hứng từ MemInsight, và cập nhật đồ thị liên kết lấy cảm hứng từ A-Mem. Quá trình tiến hóa bộ nhớ được kiểm soát nghiêm ngặt bằng điểm số đánh giá từ Evaluator (score-gated validation) nhằm tránh rò rỉ thông tin đáp án của benchmark.
    \item **Mô-đun Tiến Hóa Công Cụ (Tool Evolution Module):** Hỗ trợ AI Agent tự phát triển năng lực sử dụng công cụ tại test-time (lấy cảm hứng từ DRAFT, TTE và Alita). Agent có thể tự động Master mô tả các công cụ nguyên thủy (primitive tools), tổng hợp thành các quy trình chẩn đoán phức tạp (composite workflows), hoặc sinh mã nguồn Python (generated Python helpers) để xử lý dữ liệu. Toàn bộ các công cụ mới đều được kiểm tra cú pháp AST và chạy kiểm định trong môi trường Docker sandbox an toàn trước khi lưu trữ lâu dài.
\end{enumerate}

Hiệu quả của hệ thống được kiểm chứng thông qua các thực nghiệm chẩn đoán chi tiết trên tập con 30 sự cố chọn lọc từ benchmark NIKA. Kết quả cho thấy agent được trang bị mô-đun tiến hóa đạt độ chính xác chẩn đoán vượt trội so với baseline thông thường (điểm F1 chẩn đoán tăng từ 0.0111 lên 0.1611, cải thiện đáng kể khả năng chẩn đoán trong các sự cố mất cấu hình định tuyến và IP), đồng thời giảm số lượng bước gọi công cụ thừa và tối ưu hóa chi phí token sử dụng.

\vspace{10pt}
\noindent \textit{\textbf{Từ khóa}: Self-Evolving AI Agents, xử lý sự cố mạng, NIKA, bộ nhớ quy trình tiến hóa, tiến hóa công cụ, Kathará.}
